% ============================================================
%  凡例（frontmatter/conventions.tex）
%  与序同级：\chapter* 不编号；入目录（同序，用 part 级）；
%  由 main.tex \input 引入，紧随 preface 之后。
%
%  结构：节用 \section*，条用 fanli 列表（中文序数由 zhnumber 出）。
%  文体：白话书面语。不用赘语，不作承诺，不写理由。
%  ★ 待填：标「（待填）」者补实后删去该字样。
% ============================================================

\chapter*{凡例　Conventions}

% —— 入目录（不编号）——
\begingroup
    \renewcommand{\numberline}[1]{}   % 局部重定义，出组自动复原
    \addcontentsline{toc}{part}{凡例　Conventions}
\endgroup

% —— 页眉/页脚 ——
% 页眉标记须在内容之前设置：置于文末则本页仍用上一篇的 \markboth{序　Preface}。
\markboth{凡例　Conventions}{}
\thispagestyle{plain}

% 凡例条目表：中文序数、左缩进二字、条间略空、句首不另缩进
% 注一：label 内须用本列表自己的计数器名（fanli 的计数器为 fanlii）。
%       用通用的 enumi 取不到值，会冻结成同一个数字。
% 注二：对齐用 align=left，不可用 align=parleft——后者在标签过宽时
%       （如「十一、」「十二、」两字序号）会把标签折成单独一行。
\newlist{fanli}{enumerate}{1}
\setlist[fanli]{label=\zhnum{fanlii}、, labelsep=0pt, leftmargin=2\ccwd,
  align=left, itemsep=3pt, topsep=4pt, parsep=0pt, partopsep=0pt}

\section*{述旨}

\begin{fanli}
  \item 本册是课堂笔记的整理，供复习、查阅。与课本、课程标准不一致的地方，以课本和课程标准为准。
  \item 本篇说明本册的编辑体例。其余说明散见各章正文。
\end{fanli}

\section*{一、编次}

\begin{fanli}
  \item 按学程分卷，卷下分章，章下分节。
  \item 每一知识点按「概念—要点—例证」排列，例题随附解答。
  \item 同一知识点在多处出现时，第一次出现的地方写详细，其余各处只标「参见」。
  \item 条目顺序按讲授顺序和课本顺序，不按知识体系重排，所以各章详略不齐。
  \item 实验、计算与专题归纳各自立节。
\end{fanli}

\section*{二、标记}

\begin{fanli}
  \item 并列分条用「① ② ③」。
  \item 判断正误用两组记号：细勾 \cmark{} 用于随文判断，粗勾 \hcmark{}、粗叉 \hxmark{} 用于整行判断。两组记号取自西文字体（Source Sans），与中文字体不同源。
  \item 着重记号两种：下加横线用于必须记住的词句，加重体用于标题性词语与表头。两者不叠用，一段之内不超过三处。
  \item 颜色的作用是用三种颜色区分轻重：蓝色表示课上讲的名物与义理，红色表示课后的要诀与陷阱，灰色表示图注与附言。
  \item 三种颜色都可以不用。文意已经说清的，或只是叙述性的文字，不必着色；凡是着色，都要有理由。
  \item 蓝色用于：术语、定义、通性与规律、所设的条件、推理与计算的过程。
  \item 红色用于：结论与小结、口诀与速记、易错与易混之处、例题的答案。
  \item 灰色用于：图注，以及次要的附言，如晶胞均摊法的数据。
  \item 同一个短语只用一种颜色；两种都说得通的，用红色。
  \item 一个框内不要全部着色，全篇也不以颜色多为好。
  \item 颜色只作辅助，黑白打印时文意应当完整。
  \item 同一概念全书只用同一记号，不用异体。
  \item 序中称名、敬称和标点的写法已见序中，本篇不再重复；正文一律照此，不作别体。
\end{fanli}

\section*{三、化学语言}

\begin{fanli}
  \item 化学式、离子式、方程式依 IUPAC 标准书写。
  \item 离子电荷写在上标，位于下标之后。硫酸根写作 SO\textsubscript{4}\textsuperscript{2−}；电荷数字在前、正负号在后。电荷为一时只写符号，写作 H\textsuperscript{+}，不写 H\textsuperscript{1+}。
  \item 氧化数用罗马数字加符号，置于元素符号后的括号内，如 Fe(+III)、MnO\textsubscript{4}\textsuperscript{−} 中的 Mn(+VII)。
  \item 物质状态在化学式后用括号标注：气体 (g)、液体 (l)、固体 (s)、溶液 (aq)。热化学方程式须标状态并写出 ΔH。
  \item 结晶水用「·」连接，如 CuSO\textsubscript{4}·5H\textsubscript{2}O。
  \item 同位素的质量数与原子序数分别写在元素符号的左上与左下，如 \(\cemh{^{235}_{92}U}\)。
  \item 反应条件（加热、催化剂、通电、光照、高温高压等）写在反应箭头上下，不写在箭头后面。
  \item 化学方程式与离子方程式同时出现时，先写化学方程式，再写离子方程式，并说明两者的对应关系。
  \item 有机物结构式与简式并用：简式用于书写反应与命名，结构式用于说明空间构型与同分异构。
  \item 本册自拟的口诀、归纳用红字排出。
\end{fanli}

\section*{四、数值与单位}

\begin{fanli}
  \item 数值与单位写作 \SI{0.1}{\mole\per\litre}、\SI{25}{\degreeCelsius}，温度不写作 25°C。
  \item 单位用国际单位制，正体排；数值与单位之间空一个字的位置。
  \item 复合单位用居中圆点连接，不用斜线叠写。
  \item 同一物理量全书只用一种写法。
  \item 近似值、实测值注明条件，如温度、压强、浓度。
  \item 有效位数本册内统一：数据表按原书，例题随题目所给。（待填）
  \item 中间计算步骤不提前舍入，以免误差累积。
\end{fanli}

\section*{五、图表}

\begin{fanli}
  \item 本册插图都是自己画的，用来帮助理解；各图以示意为主，但原则上按精确比例。
  \item 需要精确比例的地方，把尺寸或数据标在图旁。
  \item 晶胞与结构图随文就近排出，说明用小号灰字排在图下。
  \item 图题写作「图 \(x.y\)」，表题写作「表 \(x.y\)」，随文编排。
  \item 正文提到图表时标出编号，如「见图 \(x.y\)」。
  \item 表内数据的出处按「十、引用」标注。
\end{fanli}

\section*{六、环境}

\begin{fanli}
  \item 知识点条目一律置框：左侧一道色边，底色浅淡；配色随条目类别而分，见下。
  \item 定义框（绿）：下术语界说。
  \item 性质框（粉）：列出物质的性质、规律与通性。
  \item 例题框（紫）：例题与随附解答。
  \item 习题框（蓝）：待做的题目；必要时可指定起始编号。
  \item 解答框（灰）：解答与证明，无编号。
  \item 按语框（橙）：补充、提醒与例外。
  \item 反例框（红）：反例与易错之处。
  \item 数据清单（bullet 列表）：逐条数据，与正文分开。
  \item 框色分类与「二、标记」讲的红、蓝、灰三色分工不同：框色区分条目的类别，三色标记条目以内的内容。
\end{fanli}

\section*{七、化学式与数值}

\begin{fanli}
  \item 化学式、离子式、方程式书写时对应用 Unicode 上下标字符。硫酸根作 SO\textsubscript{4}\textsuperscript{2−}，氢离子作 H\textsuperscript{+}。
  \item 原则上括号不分级：仅使用小括号。
  \item 反应条件写在箭头上方或下方，箭头作 \cemh{A ->[\text{上}][\text{下}] B} 之形；可逆反应用可逆号。
  \item 分条编号作 ① ② ③。
  \item 正误判断记号作细勾 \cmark{}、粗勾 \hcmark{}、粗叉 \hxmark{}，用法见「二、标记」。
\end{fanli}

\section*{八、导言所定}

\begin{fanli}
  \item 篇章：分卷、章、节、条四级。章编号作一、二、三；节号前加「§」，并带上课本章号；小节作 1.1.1 的形式。目录列到小节。卷与章显示在目录、书签和页眉中。
  \item 开本与版心：A4 双面，内侧留装订边；正文左对齐，段首缩进两字。
  \item 页眉：偶数页在左侧标章名，奇数页在右侧标节名。
  \item 页脚：内侧下方是页码；外侧下方是书名。
  \item 脚注标记是带圈序号，不是上标数字。
  \item 页码：篇首的序与凡例用罗马数字，正文从阿拉伯数字 1 开始。
  \item 引用与链接：文内引用自动显示章、节、图表名；电子版链接为红蓝两色，打印版转黑。
  \item 列表：条目之间不留空隙。
  \item 图与表：标题作「图 \(x.y\)」「表 \(x.y\)」，可以跨页续排。
  \item 字体跟随系统：正文宋体、标题黑体、图注灰色小号；上下标与数学量用西文字体。
\end{fanli}

\section*{九、出版}

\begin{fanli}
  \item 电子版保留红、蓝链接颜色。
  \item 打印版在项目根目录建立空文件 \verb|printmode.txt| 后编译，链接颜色转黑，仍可点击。
\end{fanli}

\section*{十、引用}

\begin{fanli}
  \item 引课本的地方，标出章、节、页码。
  \item 引老师口授的地方，标出学期和课题；口授没有页码，不强求标页码。
  \item 数据表（电离常数、溶度积、相对原子质量等）注明来源、温度和有效位数。
  \item 本册自己归纳的口诀和速记不标出处。
  \item 引其他书时标书名、版本、页码；引自网络的标检索日期。
  \item 引文有删节的标省略号；有改动的按「十一、校勘」处理。
\end{fanli}

\section*{十一、校勘}

\begin{fanli}
  \item 错字同时用「（）」和「〔〕」处理。补上或删去的字都写在括号内，括号外的文字不动。
  \item 补上的字如果是编者加的，必要处加按语；原文本来就有的，直接写在括号内，不加按语。
  \item 发现图文不符、数据有误或体例不一，请读者批评斧正。
  \item 已经改正的地方，改版时在原文之外并列「勘误表」，另记改正之处和依据。
\end{fanli}
